\section*{الفصل \ref{ch:b1:erosion-transport} --- التعرية والنقل}

\begin{solution}{exo:b1:erosion-transport:1}
\qty{0.4}{g.L^{-1}} يساوي \qty{0.4}{kg.m^{-3}}، فالتدفّق $500 \times 0.4 =
\qty{200}{kg.s^{-1}}$. والسنة من \num{365.25}~يوم تدوم
\qty{3.156e7}{s}، فالتدفّق السنوي \qty{6.31e9}{kg}، أو
\qty{6.31}{Mt}.
\end{solution}

\begin{solution}{exo:b1:erosion-transport:2}
$w_s = R g d^2/(18\nu) = 1.6548 \times 9.81 \times (\num{30e-6})^2/(18 \times
\num{1.0034e-6}) = \qty{8.09e-4}{m.s^{-1}}$، أي نحو \qty{0.8}{mm} في الثانية. وعدد رينولدز للحبيبة هو $w_s d/\nu = \num{8.09e-4} \times
\num{30e-6}/\num{1.0034e-6} = 0.024$، أدنى بكثير من 1: قانون ستوكس يصحّ. ويهبط عبر \qty{2}{m} في $2/\num{8.09e-4} = \qty{2472}{s}$، أي نحو
\qty{41}{min}.
\end{solution}

\begin{solution}{exo:b1:erosion-transport:3}
$\tau_b = \rho g h S = 998.2 \times 9.81 \times 0.5 \times \num{4e-4} =
\qty{1.96}{Pa}$. مع $(\rho_s - \rho) g = 1651.8 \times 9.81 =
\qty{16204}{N.m^{-3}}$: ل\qty{1}{mm}، $\theta = 1.96/(16204 \times 0.001) =
0.121$، أي نحو أربعة أضعاف عتبة \num{0.0314}، فالرمل يتحرّك؛ ول
\qty{10}{mm}، $\theta = 0.0121$، أي خُمس عتبة \num{0.0556}، فالحصى يبقى. يحمل النهر رماله فوق قاع من الحصى لا يتحرّك.
\end{solution}

\begin{solution}{exo:b1:erosion-transport:4}
(أ)~$12.6/3.9 = 3.2$: تحمل الأنهار نحو ثلاثة أضعاف في المعلّق مما في المحلول. (ب)~الخزانات الكبيرة: $0.20 \times 14 = \qty{2.8}{Gt}$ في السنة؛ وكلها: $0.26 \times 14 = \qty{3.64}{Gt}$ في السنة. (ج)~التدفّق المذاب هو \qty{3.9e15}{g} في السنة عند \qty{120}{g.m^{-3}}، فالماء الذي يحمله هو
$\num{3.9e15}/120 = \qty{3.25e13}{m^3}$، نحو \qty{32500}{km^3} في السنة. \omterm{def:b1:erosion-transport:load}{والحمولة المعلّقة} في ذلك الماء هي $\qty{1.26e16}{g}/\qty{3.25e16}{L} =
\qty{0.39}{g.L^{-1}}$.
\end{solution}

\begin{solution}{exo:b1:erosion-transport:5}
(أ)~للمقطع المساحة $10 \times 2 = \qty{20}{m^2}$ والمحيط المبلول $10 + 2 \times 2 = \qty{14}{m}$، ف$R_h = 20/14 = \qty{1.43}{m}$. (ب)~$\tau_b = 998.2 \times 9.81 \times 1.43 \times 0.001 = \qty{14.0}{Pa}$، مقابل $\qty{19.6}{Pa}$ لقناة عريضة جدًّا عمقها \qty{2}{m}: النسبة $R_h/h = 0.714$، إذ تأخذ الضفاف الباقي من الاحتكاك. (ج)~لقناة مستطيلة عرضها $b$، $R_h = bh/(b + 2h) = h/(1 + 2h/b)$، الذي يتّجه إلى $h$ حين يكبر $b$.
\end{solution}

\begin{solution}{exo:b1:erosion-transport:6}
(أ)~$d_1 = \left(18\nu^2/(Rg)\right)^{1/3} = (18/(Rg))^{1/3}\,\nu^{2/3}$، و$R$
و$g$ لا يتعلّقان بالزوجة. (ب)~$\nu = \num{1.3059e-3}/999.7 =
\qty{1.306e-6}{m^2.s^{-1}}$ عند \qty{10}{\celsius}، مقابل
\qty{1.0034e-6}{m^2.s^{-1}} عند \qty{20}{\celsius}: النسبة \num{1.302}،
و$1.302^{2/3} = 1.192$، إذن $d_1 = 0.104 \times 1.192 = \qty{0.124}{mm}$.
(ج)~في مدى ستوكس تتناسب $w_s$ عكسيًّا مع $\nu$: فالحبيبة تهبط أبطأ بمقدار $1.302$ مرة، أي نحو \qty{30}{\percent}. فالماء البارد يُبقي الرسوب الناعم معلّقًا زمنًا أطول: والنهر في الشتاء يحمل طميه أبعد قليلًا عند التصريف نفسه.
\end{solution}

\begin{solution}{exo:b1:erosion-transport:7}
مع $R g/(18\nu) = \qty{8.99e5}{m^{-1}.s^{-1}}$ والقانون المجمّع:
\begin{center}
\begin{tabular}{@{}lccc@{}}
\hline
$d$ & ستوكس (\si{m.s^{-1}}) & المجمّع (\si{m.s^{-1}}) & النسبة\\
\hline
\qty{0.1}{mm} & \num{8.99e-3} & \num{7.53e-3} & \num{1.19}\\
\qty{0.5}{mm} & \num{0.225} & \num{0.0711} & \num{3.16}\\
\qty{2}{mm} & \num{3.60} & \num{0.197} & \num{18.3}\\
\hline
\end{tabular}
\end{center}
تكبر النسبة مع الحجم، لأن عدد رينولدز للحبيبة يكبر: فقانون ستوكس، الصالح لأعداد رينولدز الصغيرة وحدها، يزيد ك$d^2$، بينما يبطئ سحب الذيل الدوّامي الحبيبات الخشنة إلى سرعة تكبر فقط ك$\sqrt d$.
\end{solution}

\begin{solution}{exo:b1:erosion-transport:8}
(أ)~$(Rg/\nu^2)^{1/3} = \qty{25263}{m^{-1}}$. عند \qty{0.2}{mm}، $D_* = 5.05$
و$\theta_c = 0.3/7.06 + 0.055(1 - e^{-0.101}) = 0.0425 + 0.0053 =
0.0478$؛ $\tau_c = 0.0478 \times 16204 \times \num{2e-4} = \qty{0.155}{Pa}$.
عند \qty{20}{mm}، $D_* = 505$ و$\theta_c = 0.3/607 + 0.055(1 - e^{-10.1}) =
0.0005 + 0.0550 = 0.0555$؛ $\tau_c = 0.0555 \times 16204 \times 0.02 =
\qty{18.0}{Pa}$. (ب)~$h = \tau_c/(\rho g S)$ مع $\rho g S =
\qty{19.58}{Pa.m^{-1}}$: \qty{7.9}{mm} من الماء تحرّك الرمل الناعم، و\qty{0.92}{m}
الحصى.
\end{solution}

\begin{solution}{exo:b1:erosion-transport:9}
يفقد الحوض \qty{0.8}{Mt} = \qty{8e5}{t} في السنة، أو $\num{8e5}/2000 =
\qty{400}{t.km^{-2}.yr^{-1}}$. وبحسب \cref{prop:b1:erosion-transport:denudation}،
$D = \qty{8e8}{kg}/(2650 \times \qty{2e9}{m^2}) = \qty{1.51e-4}{m}$ في السنة، أو
\qty{151}{m/Myr}. وذلك نحو ثلاثة أضعاف منتصف الأحواض المقاسة بالبريليوم-10، \qty{54}{m/Myr}، وإن كان دون متوسطها \qty{218}{m/Myr}: حوض سريع إلى حدّ ما، كما في بلد شديد الانحدار أو نشط تكتونيًّا.
\end{solution}

\begin{solution}{exo:b1:erosion-transport:10}
(أ)~$u_* = \sqrt{g h S} = \sqrt{9.81 \times 3 \times \num{5e-4}} =
\qty{0.1213}{m.s^{-1}}$ و$\tau_b = \rho u_*^2 = \qty{14.7}{Pa}$.
(ب)~يحتاج التعليق إلى $w_s < u_* = \qty{0.1213}{m.s^{-1}}$. فالقانون المجمّع
يعطي \qty{0.126}{m.s^{-1}} عند \qty{1}{mm}؛ وبحلّ
$w_s(d) = 0.1213$ بالتجربة يعطي $d = \qty{0.94}{mm}$. (ج)~للحصى
$\theta_c \approx 0.0555$، ف$\tau_c = \qty{14.7}{Pa}$ عند $d \approx
14.7/(0.0555 \times 16204) = \qty{16.3}{mm}$؛ وعند ذلك الحجم يعطي منحنى العتبة
$\theta_c = 0.0555$ فعلًا. (د)~الحبيبات الأخشن من نحو \qty{16}{mm} تبقى على
القاع وتكوّن أرضية حصى مستقرّة؛ والحصى الأدقّ والرمل الأخشن من نحو مليمتر
يسافران كحمولة قاعية؛ والرمل الأدقّ والطمي والطين يسافرون معلّقين.
\end{solution}

\begin{solution}{exo:b1:erosion-transport:11}
(أ)~$0.30 \times \qty{1000}{m} = \qty{300}{m}$ من الصخر، متوسطًا على المنطقة. (ب)~الارتداد هو $300 \times 2700/3300 = \qty{245}{m}$. (ج)~السطح المتوسط لا ينخفض إلا بمقدار $300 - 245 = \qty{55}{m}$. (د)~القمم، التي لم تتآكل، تُرفع مع المنطقة كلها: فهي الآن أعلى ب\qty{245}{m} مما كانت. وحفر الوديان يستطيع أن يرفع قمم سلسلة.
\end{solution}

\begin{solution}{exo:b1:erosion-transport:12}
(أ)~$0.9 \times \qty{5e8}{kg}/\qty{1300}{kg.m^{-3}} = \qty{346000}{m^3}$ في السنة،
أو \qty{0.69}{\percent} من الخزان. (ب)~نصف الخزان، \qty{25e6}{m^3}، يمتلئ في $\num{25e6}/\num{346000} = \num{72}\text{ سنة}$ بهذا المعدّل. (ج)~مع امتلاء الخزان ينقص حجمه، وبالتالي زمن بقاء الماء، ويمرّ الحبيبات الأدقّ: فتنخفض كفاءة حبسه. وفي الوقت نفسه تُدفن الرواسب القديمة وتنضغط، فتطرد الماء، فترتفع كثافتها الكلّية ويأخذ كل طن مكانًا أقلّ. وكلا التأثيرين يبطئ فقدان الحجم مع شيخوخة الخزان.
\end{solution}

\begin{solution}{pb:b1:erosion-transport:1}
\textbf{1.} سنة من الماء هي $1.2 \times \num{3.156e7} = \qty{3.79e7}{m^3}$؛ موزّعة على \qty{8e7}{m^2}، طبقة سمكها \qty{0.473}{m}: جريان قدره \qty{473}{mm} في السنة.

\textbf{2.} $1.2 \times \qty{0.150}{kg.m^{-3}} = \qty{0.18}{kg.s^{-1}}$، أو
\qty{5680}{t} في السنة.

\textbf{3.} $1.2 \times 0.080 = \qty{0.096}{kg.s^{-1}}$، أو \qty{3030}{t} في
السنة.

\textbf{4.} \omterm{def:b1:erosion-transport:load}{الحمولة القاعية} هي \qty{568}{t} في السنة، والمجموع الكلّي $5680 + 3030 +
568 = \qty{9278}{t}$ في السنة.

\textbf{5.} يكون تركيز الرسوب المعلّق مرتفعًا خلال الأيام القليلة من الفيضان، حين يكون التصريف مرتفعًا أيضًا، ومنخفضًا في بقية الوقت. ومتوسط زمني يعطي كل يوم الوزن نفسه، فيغلب عليه الأيام الكثيرة ذات التصريف المنخفض والتركيز المنخفض. وبحسب \cref{prop:b1:erosion-transport:flux}، الكتلة السنوية هي $\int Q\,C\,\mathrm dt$؛ والمتوسط المرجّح بالجريان هو هذا التكامل مقسومًا على $\int Q\,\mathrm dt$، فهو المتوسط المرجّح بالجريان الذي يجب أن يُضرب في متوسط التصريف.

\textbf{6.} $3030/9278 = \num{0.33}$، ثلثه. وبالنسبة لأنهار العالم، $3.9/(12.6 + 3.9) = \num{0.24}$، وأقلّ لو حُسبت \omterm{def:b1:erosion-transport:load}{الحمولة القاعية}. وحصّة الحوض أعلى: تلال منخفضة، حيث يتيح التآكل الميكانيكي البطيء للصخر وقتًا ليذوب، وحجر جيري وطباشير يذوبان كليًّا.

\textbf{7.} $9278/80 = \qty{116}{t.km^{-2}.yr^{-1}}$.

\textbf{8.} $D = \qty{9.278e6}{kg}/(2700 \times \qty{8e7}{m^2}) =
\qty{4.30e-5}{m}$ في السنة، \qty{43.0}{m/Myr}.

\textbf{9.} $43.0/54 = 0.80$: أدنى قليلًا من منتصف الأحواض المقاسة، وأدنى بخمس مرات من متوسطها. حوض بطيء إلى حدّ ما، كما هو متوقّع لتلال منخفضة في منطقة هادئة.

\textbf{10.} يتوسّط القياس عبر السنوات القليلة للقياس، ويتوقّف على الفيضانات التي حدثت فيها. ويتوسّط البريليوم-10 على الزمن الذي يستغرقه تآكل الأمتار القليلة التي يتركّز فيها توليده: وعند \qty{43}{m/Myr}، يستغرق كل متر \num{23000} سنة، فيغطي المتوسط عشرات الآلاف من السنين. ويمكن أن يختلفا بسبب الفيضانات الكبيرة النادرة، الغائبة عن سجل قصير، ولأن الزراعة وإزالة الغابات رفعتا التآكل في القرون الأخيرة، ولأن الرسوب يُخزَّن في السهل الفيضي ويُستعاد منه.

\textbf{11.} جزء من \omterm{def:b1:erosion-transport:load}{الحمولة المذابة} لا يأتي من الصخر: فكثير من البيكربونات يحمل كربونًا مأخوذًا من الهواء أو التربة، والمطر يجلب ملح البحر. وعدّها سيمنح الحوض صخرًا لم يفقده.

\textbf{12.} $\num{1.5e6}/1.2 = \qty{1.25e6}{s}$، \num{14.5} يومًا.

\textbf{13.} يعطي القانون المجمّع \qty{7.53e-3}{m.s^{-1}}: تسقط الحبيبة \qty{8}{m} في \qty{1062}{s}، نحو \qty{18}{min}. فهي تُحبَس.

\textbf{14.} $w_s = \num{8.99e5} \times (\num{2e-5})^2 =
\qty{3.60e-4}{m.s^{-1}}$: \qty{8}{m} في \qty{22200}{s}، نحو \num{6.2}~ساعات. تُحبَس.

\textbf{15.} $w_s = \num{8.99e5} \times (\num{1e-6})^2 =
\qty{8.99e-7}{m.s^{-1}}$: \qty{8}{m} تأخذ \qty{8.9e6}{s}، \num{103}~يومًا، أي سبعة أضعاف زمن البقاء. فمعظمه يعبر السدّ مع الماء.

\textbf{16.} يجب أن تهبط الحبيبة \qty{8}{m} في \qty{1.25e6}{s}، بسرعة $\qty{6.4e-6}{m.s^{-1}}$ على الأقل. وبقانون ستوكس $d = \sqrt{18\nu
w_s/(Rg)} = \sqrt{\num{6.4e-6}/\num{8.99e5}} = \qty{2.7}{\micro m}$: طين ناعم. والحبيبات الأخشن تُحبَس أينما دخلت.

\textbf{17.} $568 + 0.7 \times 5680 = \qty{4544}{t}$ في السنة، وهذا يملأ
$\num{4.544e6}/1300 = \qty{3496}{m^3}$ في السنة، \qty{0.233}{\percent} من
الخزان. ونصفه، \qty{7.5e5}{m^3}، يمتلئ في نحو \num{215} سنة.

\textbf{18.} $100/43.0 = \qty{2.33}{Myr}$.

\textbf{19.} $1 - 2700/3300 = 0.18$: ينخفض السطح بنسبة \qty{18}{\percent} فقط من السمك المُزال.

\textbf{20.} $100/0.1818 = \qty{550}{m}$ من الصخر.

\textbf{21.} ترتفع المنطقة كلها بمقدار $550 \times 2700/3300 = \qty{450}{m}$
من الارتداد. وتُرفع القمم التي بالكاد تتآكل معها، بمقدار يصل إلى \qty{450}{m}، بينما تتعمّق الوديان: فالتضاريس تزداد قبل أن تنحسر.

\textbf{22.} يتوقّف \omterm{def:b1:erosion-transport:denudation}{معدّل التعرية} قبل كل شيء على المنحدرات. فمع خفض التلال وتخفيف منحدراتها سينخفض المعدّل، ما لم يرفعها الارتفاع أو يغيّر المناخ فتعود إلى الانحدار: والزمن المحسوب بمعدّل اليوم حدٌّ أدنى.

\textbf{23.} $550/42.95 = 12.80$. بمعدّل اليوم وبالارتداد الإيزوستاتي، يستغرق خفض تلال فاردن بمقدار \qty{100}{m} \textbf{12.8}~مليون سنة.
\end{solution}
